% GENERATED FILE. Edit canonical lessons, then run npm run generate:books.
% canonical-source-hash: fnv1a64:8c7ff0d6b2a6f58a
% canonical-lessons: SA-C113-sasakah, SA-C113-gardabhah, SA-C113-varahah, SA-C113-vrsabhah, SA-C113-mahisah

\chapter{Rabbit, Donkey, Boar, Bull, Buffalo}
\label{ch:sa-rabbit-donkey-boar-bull-buffalo}
\hlchaptermodality{113}

\begin{chapteropening}
\textbf{By the end of this chapter:} \emph{I can say a rabbit, a donkey, a boar, a bull and a buffalo.}

Complete the last lesson of chapter 113: i can say a rabbit, a donkey, a boar, a bull and a buffalo.
\end{chapteropening}

\section[\texorpdfstring{\sk{शशकः} (śaśakaḥ)}{शशकः (śaśakaḥ)}]{\sk{शशकः} (śaśakaḥ) — a rabbit}

\label{lesson:SA-C113-sasakah}



\begin{quote}
Before the new one: say the Sanskrit for a bear, then the Sanskrit for a jackal.
\end{quote}

\subsection*{You'll want to know: \sk{शशकः}}
\textbf{\sk{शशकः}} — \emph{śaśakaḥ} — \textquotedblleft{}a rabbit\textquotedblright{}.

\begin{grammarlens}[title={What you've built}]
One more word for this chapter.
\end{grammarlens}

\subsection*{Your turn}
\begin{itemize}
\raggedright
  \item \emph{Say it:} \emph{śaśakaḥ}
  \item \emph{Say it:} \emph{śaśakaḥ}, once more
  \item \emph{Say it:} say \emph{śaśakaḥ}
  \item \emph{Recall:} say the Sanskrit for a bear, then the Sanskrit for a jackal, then say \emph{śaśakaḥ} again
\end{itemize}

\begin{tcolorbox}[breakable,colback=teal!4,colframe=teal!35!black,title={Before you move on}]
What does \sk{शशकः} mean? (\textquotedblleft{}A rabbit\textquotedblright{}.) Say it once more.
\end{tcolorbox}
\section[\texorpdfstring{\sk{गर्दभः} (gardabhaḥ)}{गर्दभः (gardabhaḥ)}]{\sk{गर्दभः} (gardabhaḥ) — a donkey}

\label{lesson:SA-C113-gardabhah}



\begin{quote}
Before the new one: say the Sanskrit for a jackal, then the Sanskrit for a rabbit.
\end{quote}

\subsection*{You'll want to know: \sk{गर्दभः}}
\textbf{\sk{गर्दभः}} — \emph{gardabhaḥ} — \textquotedblleft{}a donkey\textquotedblright{}.

\begin{grammarlens}[title={What you've built}]
One more word for this chapter.
\end{grammarlens}

\subsection*{Your turn}
\begin{itemize}
\raggedright
  \item \emph{Say it:} \emph{gardabhaḥ}
  \item \emph{Say it:} \emph{gardabhaḥ}, once more
  \item \emph{Say it:} say \emph{gardabhaḥ}
  \item \emph{Recall:} say the Sanskrit for a jackal, then the Sanskrit for a rabbit, then say \emph{gardabhaḥ} again
\end{itemize}

\begin{tcolorbox}[breakable,colback=teal!4,colframe=teal!35!black,title={Before you move on}]
What does \sk{गर्दभः} mean? (\textquotedblleft{}A donkey\textquotedblright{}.) Say it once more.
\end{tcolorbox}
\section[\texorpdfstring{\sk{वराहः} (varāhaḥ)}{वराहः (varāhaḥ)}]{\sk{वराहः} (varāhaḥ) — a boar}

\label{lesson:SA-C113-varahah}



\begin{quote}
Before the new one: say the Sanskrit for a rabbit, then the Sanskrit for a donkey.
\end{quote}

\subsection*{You'll want to know: \sk{वराहः}}
\textbf{\sk{वराहः}} — \emph{varāhaḥ} — \textquotedblleft{}a boar\textquotedblright{}.

\begin{grammarlens}[title={What you've built}]
One more word for this chapter.
\end{grammarlens}

\subsection*{Your turn}
\begin{itemize}
\raggedright
  \item \emph{Say it:} \emph{varāhaḥ}
  \item \emph{Say it:} \emph{varāhaḥ}, once more
  \item \emph{Say it:} say \emph{varāhaḥ}
  \item \emph{Recall:} say the Sanskrit for a rabbit, then the Sanskrit for a donkey, then say \emph{varāhaḥ} again
\end{itemize}

\begin{tcolorbox}[breakable,colback=teal!4,colframe=teal!35!black,title={Before you move on}]
What does \sk{वराहः} mean? (\textquotedblleft{}A boar\textquotedblright{}.) Say it once more.
\end{tcolorbox}
\section[\texorpdfstring{\sk{वृषभः} (vṛṣabhaḥ)}{वृषभः (vṛṣabhaḥ)}]{\sk{वृषभः} (vṛṣabhaḥ) — a bull}

\label{lesson:SA-C113-vrsabhah}



\begin{quote}
Before the new one: say the Sanskrit for a donkey, then the Sanskrit for a boar.
\end{quote}

\subsection*{You'll want to know: \sk{वृषभः}}
\textbf{\sk{वृषभः}} — \emph{vṛṣabhaḥ} — \textquotedblleft{}a bull\textquotedblright{}.

\begin{grammarlens}[title={What you've built}]
One more word for this chapter.
\end{grammarlens}

\subsection*{Your turn}
\begin{itemize}
\raggedright
  \item \emph{Say it:} \emph{vṛṣabhaḥ}
  \item \emph{Say it:} \emph{vṛṣabhaḥ}, once more
  \item \emph{Say it:} say \emph{vṛṣabhaḥ}
  \item \emph{Recall:} say the Sanskrit for a donkey, then the Sanskrit for a boar, then say \emph{vṛṣabhaḥ} again
\end{itemize}

\begin{tcolorbox}[breakable,colback=teal!4,colframe=teal!35!black,title={Before you move on}]
What does \sk{वृषभः} mean? (\textquotedblleft{}A bull\textquotedblright{}.) Say it once more.
\end{tcolorbox}
\section[\texorpdfstring{\sk{महिषः} (mahiṣaḥ)}{महिषः (mahiṣaḥ)}]{\sk{महिषः} (mahiṣaḥ) — a buffalo}

\label{lesson:SA-C113-mahisah}



\begin{quote}
Before the new one: say the Sanskrit for a boar, then the Sanskrit for a bull.
\end{quote}

\subsection*{You'll want to know: \sk{महिषः}}
\textbf{\sk{महिषः}} — \emph{mahiṣaḥ} — \textquotedblleft{}a buffalo\textquotedblright{}.

\begin{grammarlens}[title={What you've built}]
One more word for this chapter.
\end{grammarlens}

\subsection*{Your turn}
\begin{itemize}
\raggedright
  \item \emph{Say it:} \emph{mahiṣaḥ}
  \item \emph{Say it:} \emph{mahiṣaḥ}, once more
  \item \emph{Say it:} say \emph{mahiṣaḥ}
  \item \emph{Recall:} say the Sanskrit for a boar, then the Sanskrit for a bull, then say \emph{mahiṣaḥ} again
\end{itemize}

\begin{tcolorbox}[breakable,colback=teal!4,colframe=teal!35!black,title={Before you move on}]
What does \sk{महिषः} mean? (\textquotedblleft{}A buffalo\textquotedblright{}.) Say it once more.
\end{tcolorbox}
